% ECONOMICS_PROBLEM: {"id": "C72-131", "title": "Expected-polynomial exact Nash computation for random rational games", "jel_code": "C72", "source_id": "OP-0137"}
% !TeX program = xelatex
\documentclass[11pt,a4paper]{article}
\usepackage[margin=23mm]{geometry}
\usepackage{fontspec}
\setmainfont{Latin Modern Roman}
\usepackage{amsmath,amssymb,mathtools}
\usepackage{xcolor,enumitem,booktabs,longtable,array,xurl,hyperref}
\definecolor{muted}{HTML}{53636C}
\hypersetup{colorlinks=true,urlcolor=blue,linkcolor=blue}
\setlength{\parindent}{0pt}
\setlength{\parskip}{5pt}
\newcommand{\R}{\mathbb R}
\newcommand{\E}{\mathbb E}
\newcommand{\N}{\mathbb N}
\newcommand{\ind}{\mathbf 1}
\newcommand{\Var}{\operatorname{Var}}
\newcommand{\Cov}{\operatorname{Cov}}
\newcommand{\supp}{\operatorname{supp}}
\DeclareMathOperator*{\argmin}{arg\,min}
\DeclareMathOperator*{\argmax}{arg\,max}
\newcommand{\entrymeta}[3]{{\sffamily\small\color{muted}\textbf{\texttt{#1}}\quad|\quad #2\par #3\par}}
\newcommand{\Problemref}[1]{\texttt{#1}}
\newcommand{\JELref}[2]{\texttt{#2} (JEL \texttt{#1})}
\begin{document}
\section*{Expected-polynomial exact Nash computation for random rational games}
\textbf{Problem ID:} \texttt{C72-131}\quad \textbf{JEL:} \texttt{C72}\par
% BEGIN ECONOMICS STATEMENT
\label{problem:OP-0137}
% GAME_THEORY_PROBLEM: {"local_id": "GT-007", "original_number": 7, "kind": "R", "entry_kind": "statement_only", "published": false, "review_required": true, "category": "game_theory"}
\label{game:GT-007}
{\sffamily\small\color{muted}
\textbf{GT-007} | Average-case equilibrium complexity\\
\textbf{R} | Formal distribution-dependent version of the random-bimatrix agenda.
\par}
\subsection*{Definitions and mathematical statement}
Fix two probability distributions $\nu_A,\nu_B$ with finite rational supports in $[0,1]$ and rational probabilities. For each $n$, sample all entries of $A\in\mathbb Q^{n\times n}$ independently from $\nu_A$, all entries of $B$ independently from $\nu_B$, and the two matrices independently. The supports and probabilities are fixed rather than growing with $n$. A Las Vegas algorithm must terminate almost surely and output a rational exact Nash equilibrium on every rational bimatrix input; its randomness is denoted $\xi$.
The model-specific question is whether
\[
 \exists\mathcal A\ \exists C,d<\infty\ \forall n\geq1,\qquad
 \mathbb E_{A,B,\xi}[T_{\mathcal A}(A,B;\xi)]\leq Cn^d.
\]
The broader research family asks which pairs $(\nu_A,\nu_B)$ satisfy this property, and whether every such fixed finite-support pair does.
\subsection*{Short English statement}
Can exact equilibrium computation be efficient on average even when worst-case games are hard?
\subsection*{Variants and scope boundaries}
Continuous real-valued distributions require a real-arithmetic or finite-precision input model before ``exact'' bit complexity is meaningful. Expected polynomial time is stronger than polynomial time with high probability: rare costly inputs count in the expectation. Density sequences $p=p(n)$ in Bernoulli games are treated separately in GT-008.
\subsection*{Source relationship and status}
[S1] calls the general random-bimatrix problem longstanding and studies Bernoulli payoffs. It does not pose the displayed all-finite-support-distributions assertion verbatim; the finite rational law is an editorial encoding repair of the original catalog's unspecified ``natural random model.''
\subsection*{References}
\begingroup\footnotesize\raggedright
{\footnotesize\raggedright
\textbf{[S1]} Andrea Collevecchio, G\'abor Lugosi, Adrian Vetta, and Rui-Ray Zhang (2025). \emph{Finding a Nash equilibrium of a random win-lose game in expected polynomial time}. arXiv:2510.12846. Locator: Abstract and Section 1.1; Section 2.2 for an always-correct algorithm with a fallback. \url{https://arxiv.org/pdf/2510.12846}\par}
\par\endgroup

\subsection*{Common conventions: Source mathematical conventions}

\textbf{Finite games.} For a finite set $A$, $\Delta(A)$ is its probability
simplex. Players choose independent mixed strategies in Nash equilibrium;
correlation is allowed only when explicitly part of the solution concept.
In a game with payoff functions $u_i$ and strategy profile $x$, an additive
$\varepsilon$ Nash equilibrium satisfies
\[
 u_i(x)\geq u_i(x_i',x_{-i})-\varepsilon
 \quad\text{for every player $i$ and every admissible deviation $x_i'$}.
\]
For numerical approximation, payoffs are normalized as locally specified.
An $\varepsilon$ payoff residual is distinct from distance at most
$\varepsilon$ to the set of exact equilibria. Point convergence, convergence
of empirical play, and convergence in distance to an equilibrium set are
also distinct.

\textbf{Computational representations.} Unless stated otherwise, a complexity
question takes rational data with binary encoding; polynomial time is measured
in total bit length. A fixed-accuracy polynomial algorithm is not necessarily
polynomial in $1/\varepsilon$, and neither is necessarily polynomial in
$\log(1/\varepsilon)$. Succinct games and value, demand, or sampling oracles
must declare their interfaces. Exact real solutions require an explicit
representation; an unspecified list of arbitrary real numbers is not a
finite computational output. A claimed algorithmic barrier is meaningful
only for its specified model and reductions.

\textbf{Dynamic games.} Histories, information, observation, and allowable
randomization are part of the model. A stationary strategy depends only on
the current state; a positional strategy in a deterministic graph game
chooses one action at each controlled vertex. Nash equilibrium at an initial
state and equilibrium after every history are different requirements.
An expectation of a pathwise $\liminf$ payoff, a $\liminf$ of expected averages,
a limiting expected average, and a uniform finite-horizon equilibrium are
different criteria. None is silently substituted for another.

\textbf{Allocation and mechanisms.} A complete allocation assigns every item
exactly once unless otherwise stated. Zero-valued goods, negative values,
capacity constraints, feasibility, lotteries, and copies of goods affect
fairness definitions and are specified locally. Dominant-strategy incentive
compatibility, Bayesian incentive compatibility, universal truthfulness,
and truthfulness in expectation impose different inequalities. Whenever
an objective ratio has zero denominator, its convention is stated or those
instances are excluded.

\textbf{Infinite-dimensional models.} Expectations, integrals, stochastic
equations, and optimization objectives require their local regularity,
integrability, filtrations, and admissible domains. A backward--forward
mean-field system is a boundary-value problem; it is not an initial-value
semiflow without an additional construction. Long-time, large-population,
and vanishing-noise limits require an order of limits and a convergence
topology. If a minimum need not be attained, the objective is its infimum
or supremum and the approximation requirement is stated separately.
% END ECONOMICS STATEMENT
\end{document}
